\scanpage{37}
On peut songer à le rendre additif, en le définissant d'abord sur les $E_{\unsure{\eta|y}}$ simples (voire \unsure{abst.|abs\textsuperscript{t}.|abs.} simples, quitte à se permettre des extensions finies sur $\unsure{\mathbb{Q}_{\ell}|\Phi_{\ell}}$) et en prolongeant par linéarité, ce qui revient à définir des facteurs locaux d'un $L_{E_\eta}(\unsure{t|T})$ par $L_{E_\eta(x)}(t)$, en posant $E_\eta(x)=\sum E_\eta^{i*}(x)$, les $E_\eta^i$ étant les facteurs simples \unsure{dans|de} une suite de composition de $E$. Ceci n'est pas encore satisfaisant, car les $E_\eta(x)$ ainsi construits ne sont pas de caractère local en $(E,x)$, --- i.e. ne commutent pas à l'extension \unsure{étale} \deleted{\unsure{sur|en} $X$} (à ext. résiduelle \inserted{en $x$} \deleted{\unsure{en}} triviale) \unsure{sur|en} $X$.
% archive sheet number: 36
